\section{Conclusion}

We presented an exploratory intraoperative pulmonary nodule localization framework based on tactile sensing and deep learning, organized around a detection-first logic. The results show that flexible tactile spatiotemporal sensing can support reliable nodule detection, can further support size inversion, and under explicit hierarchical organization, can enable coarse depth discrimination. Mechanism-guided analysis and XGBoost first established that the tactile data contained learnable information; raw-input hierarchical neural models then showed that the original tensors themselves could support learning; interpretability analyses further indicated that the neural representations encoded part of the physically meaningful structure; and a deployment-oriented unified inverter improved predicted-route robustness for system integration. Taken together, these results support the feasibility of a tactile sensing route for intraoperative pulmonary nodule localization while defining a clear boundary: this work should be read as a validated engineering and system-organization pathway rather than a finalized clinical system.
